  \chapter { Layout }

  \section{ Visual }

  Empece a transcribir las notas de lógica y note que el layout del libro es muy pobre,
  es decir, renderiza bien pero dado que visualmente esta mas cargado que un PDF de
  LaTeX se ve mal. Es decir, los bloques o cajas se expanden y quizá, una tabla que
  solo ocupa un cuarto de pagina se visualiza abajo, centrada. En un PDF esto no se
  ve mal pero en un webbook cuya intención es extender un poco las funciones de LaTeX,
  queda algo incomodo de ver.

  Para ello, surge la necesidad de integrar minipages. Intencionalmente se creo un
  alias, \backslash join, este lee dos bloques de argumentos, pero esto implica
  crear mas entornos, comandos nuevos lo cual hace que el escritor deba memorizar
  mas cosas, es por ello que priorisando LaTeX puro o vanilla es que decidí integrar
  el entorno \texttt{minipage}. Y con ello, al finalizar el arbol, un proceso que
  lee cada nodo y si encuentra dos minipages juntas, las envuelve en un wrapper
  join, que es directamente mas sencillo de leer desde JS.

  \begin{minipage}[t]{0.48\textwidth}
    \begin{axiom}[title=Doble negación]
      $\lnot\lnot p \equiv p$
    \end{axiom}
  \end{minipage}\hfill

  \begin{minipage}[t]{0.48\textwidth}
    \begin{tabular}{cc}
      Head & table \\
      \hline
      Una & tabla \\
      que & se pega
    \end{tabular}
  \end{minipage}

  Este es un ejemplo, aun que no habia contemplado el layout, creo que seria importante
  no solo parsear latex sino permitir que el webbook pueda ofrecer mas flexibilidad
  al momento de ordenar los elementos.

